% -----------------------------------------------------------------------------
% Document class
% scrartcl = clean article class for reports; bibliography in TOC, good tables
% -----------------------------------------------------------------------------
\documentclass[
  bibliography=totoc,
  captions=tableheading,
]{scrartcl}

% Fix minor KOMA compatibility issues
\usepackage{scrhack}

% Warn if recompilation needed (refs, TOC, etc.)
\usepackage[aux]{rerunfilecheck}

% Core math packages (align, symbols, extensions)
\usepackage{amsmath}
\usepackage{amssymb}
\usepackage{mathtools}

% Modern font handling (requires lualatex/xelatex)
\usepackage{fontspec}
\recalctypearea{} % recompute layout after font setup

% Language settings (hyphenation, captions)
\usepackage[english]{babel}

% Modern math fonts + ISO style (clean scientific notation)
\usepackage[
  math-style=ISO,
  bold-style=ISO,
  sans-style=italic,
  nabla=upright,
  partial=upright,
  mathrm=sym,
]{unicode-math}

\setmathfont{Latin Modern Math} % default math font
\setmathfont{XITS Math}[range={scr, bfscr}] % nicer script letters
\setmathfont{XITS Math}[range={cal, bfcal}, StylisticSet=1]

% Numbers & units formatting (\num, \SI, etc.)
\usepackage[
  locale=US,
  separate-uncertainty=true,
  per-mode=symbol-or-fraction,
  detect-all,
]{siunitx}

% Quotes (recommended with biblatex)
\usepackage[english=american]{csquotes}

% Float placement control (figures/tables positioning)
\usepackage{float}
\floatplacement{figure}{htbp}
\floatplacement{table}{htbp}

% Keep figures inside sections (prevents drifting)
\usepackage[section,below]{placeins}

% text wrapping around figures
\usepackage{wrapfig}      

% Caption styling (bold labels, smaller font)
\usepackage[
  labelfont=bf,
  font=small,
  width=0.9\textwidth,
]{caption}

% Subfigures (compare plots side-by-side)
\usepackage{subcaption}

% Include images
\usepackage{graphicx}

% Better tables (clean formatting + number alignment)
\usepackage{tabularray}
\UseTblrLibrary{booktabs, siunitx}

% Typography improvements (spacing, kerning)
\usepackage{microtype}

% Bibliography system (modern replacement for BibTeX)
\usepackage[backend=biber]{biblatex}
\addbibresource{../../../latex/references.bib}

% Clickable links + PDF metadata
\usepackage[
  unicode,
  pdfusetitle,
  pdfcreator={},
  pdfproducer={},
]{hyperref}

% Better PDF bookmarks
\usepackage{bookmark}

% Custom commands (needed for definitions below)
\usepackage{expl3}
\usepackage{xparse}

% Upright e (useful for exp, softmax, etc.)
\ExplSyntaxOn
\NewDocumentCommand \E {} {\symup{e}}
\ExplSyntaxOff

% Upright differential d (clean integrals)
\ExplSyntaxOn
\NewDocumentCommand \dif {m} {\mathinner{\symup{d} #1}}
\ExplSyntaxOff

% Vector command (bold vectors in math mode)
\ExplSyntaxOn
\let\vaccent=\v
\RenewDocumentCommand \v {}{
  \TextOrMath{\vaccent}{\symbf}
}
\ExplSyntaxOff

% Better spacing for scalable brackets
\usepackage{mleftright}

% Shortcuts for scalable brackets: \l( ... \r)
\ExplSyntaxOn
\let\ltext=\l
\RenewDocumentCommand \l {}{\TextOrMath{\ltext}{\mleft}}
\let\raccent=\r
\RenewDocumentCommand \r {}{\TextOrMath{\raccent}{\mright}}
\ExplSyntaxOff

% Absolute value and norm (auto-scaling with *)
\DeclarePairedDelimiter{\abs}{\lvert}{\rvert}
\DeclarePairedDelimiter{\norm}{\lVert}{\rVert}

% Definition equality symbols
\newcommand{\defeq}{\vcentcolon=}
\newcommand{\eqdef}{=\vcentcolon}

% No paragraph indent (clean report style)
\setlength{\parindent}{0pt}

\title{Task 00: Example Workflow}
\author{Akram Aki}
\date{\today}

\begin{document}

\maketitle

\section{Goal}

This task demonstrates the full workflow of the repository:
code $\rightarrow$ plots $\rightarrow$ report.

The goal is to show how a weekly task can be organized, how figures are generated
with Python, how they are included in \LaTeX{}, and where the final PDF is saved.

\section{Generating Plots}

Plots are created using Python and saved as PDF files in
Most of it should be done in a notebook for easy testing and iteration, but the final plotting code could be moved to a Python script for reproducibility.
This can be done using the makefiles.
\begin{center}
  \texttt{../build/}
\end{center}

The plotting style is controlled by the shared matplotlib configuration file

\begin{center}
  \texttt{latex/matplotlibrc}
\end{center}

This keeps fonts, figure sizes, and general plot formatting consistent across all
tasks.

\subsection{Example Function}

\begin{figure}
  \centering
  \includegraphics[width=0.8\textwidth]{../build/sine_plot.pdf}
  \caption{Sine plot generated with Python.}
  \label{fig:sine-plot}
\end{figure}

Figure~\ref{fig:sine-plot} shows a simple example plot generated from
\texttt{code/make\_plots.py}.

\subsection{Training Curve}

\begin{figure}
  \centering
  \includegraphics[width=0.8\textwidth]{../build/training_curve.pdf}
  \caption{Example training loss curve.}
  \label{fig:training-curve}
\end{figure}

Figure~\ref{fig:training-curve} shows an example of a plot that could represent a
training loss curve.

\section{Bibliography}

References are stored globally in

\begin{center}
  \texttt{latex/references.bib}
\end{center}

They are loaded through the shared header file. This means every task can use the
same bibliography database.

For example, a source can be cited with

\begin{center}
  \texttt{\textbackslash cite\{example-key\}}
\end{center}

or, depending on the context,

\begin{center}
  \texttt{\textbackslash autocite\{example-key\}}
\end{center}

The bibliography is printed at the end of the report with

\begin{center}
  \texttt{\textbackslash printbibliography}
\end{center}

\section{Build Output}

The task is built using \texttt{make}. Running

\begin{center}
  \texttt{make}
\end{center}

inside the task folder first generates the figures and then compiles the report.

Temporary build files are written to

\begin{center}
  \texttt{build/}
\end{center}

The final PDF is copied to

\begin{center}
  \texttt{report/main.pdf}
\end{center}

This is the file that should be opened or submitted.

\section{How it Works}

\begin{itemize}
  \item Python scripts are stored in \texttt{code/}
  \item Plots are generated by \texttt{code/make\_plots.py}
  \item Figures are saved in \texttt{../build/}
  \item The report is written in \texttt{report/main.tex}
  \item External PDF figures are included with \texttt{\textbackslash includegraphics}
  \item Bibliography entries are stored in \texttt{latex/references.bib}
  \item Temporary files are collected in \texttt{build/}
  \item The final report is saved as \texttt{report/main.pdf}
\end{itemize}

\section{Conclusion}

This setup allows reproducible experiments and clean reports. Each task stays
self-contained, while shared settings such as the \LaTeX{} header, bibliography,
and matplotlib configuration are reused across the repository.

\printbibliography

\end{document}
